\documentclass[11pt]{article}

% Change "review" to "final" to generate the final (sometimes called camera-ready) version.
% Change to "preprint" to generate a non-anonymous version with page numbers.
\usepackage[review]{acl}

% Standard package includes
\usepackage{times}
\usepackage{latexsym}

% For proper rendering and hyphenation of words containing Latin characters (including in bib files)
\usepackage[T1]{fontenc}
% For Vietnamese characters
% \usepackage[T5]{fontenc}
% See https://www.latex-project.org/help/documentation/encguide.pdf for other character sets

% This assumes your files are encoded as UTF8
\usepackage[utf8]{inputenc}

% This is not strictly necessary, and may be commented out,
% but it will improve the layout of the manuscript,
% and will typically save some space.
\usepackage{microtype}

% This is also not strictly necessary, and may be commented out.
% However, it will improve the aesthetics of text in
% the typewriter font.
\usepackage{inconsolata}

%Including images in your LaTeX document requires adding
%additional package(s)
\usepackage{graphicx}
\usepackage{booktabs}
\usepackage{amsmath}
\usepackage{tikz}
\usetikzlibrary{arrows.meta,positioning,decorations.pathreplacing,calc}
\usepackage{enumitem}
% generated by tools/paper_numbers.py; do not edit
\newcommand{\dscrFiftyCR}{4}
\newcommand{\dscrFiftyPD}{37}
\newcommand{\dscrFiftyPR}{4}
\newcommand{\dscrFiftySD}{9}
\newcommand{\dscrMajorityCount}{35}
\newcommand{\dscrMajorityPct}{67.3}
\newcommand{\dscrVthreeCR}{4}
\newcommand{\dscrVthreePD}{35}
\newcommand{\dscrVthreePR}{4}
\newcommand{\dscrVthreeSD}{9}
\newcommand{\nPatAll}{91}
\newcommand{\nPatEligible}{55}
\newcommand{\nPatFifty}{54}
\newcommand{\nPatVthree}{52}
\newcommand{\permDscrMean}{15.0}
\newcommand{\permDscrN}{36}
\newcommand{\permDscrObserved}{14}
\newcommand{\permDscrP}{0.73}
\newcommand{\permHi}{37}
\newcommand{\permLo}{27}
\newcommand{\permMean}{32.4}
\newcommand{\permN}{95}
\newcommand{\permObserved}{32}
\newcommand{\permP}{0.64}
\newcommand{\tcmAgree}{52}
\newcommand{\tcmClassUnique}{18}
\newcommand{\tcmFollowAbsentGemmaTwelve}{11}
\newcommand{\tcmFollowAbsentGemmaTwentySeven}{21}
\newcommand{\tcmFollowAbsentMedGemma}{9}
\newcommand{\tcmFollowAbsentScout}{15}
\newcommand{\tcmFollowCorrectGemmaTwelve}{26}
\newcommand{\tcmFollowCorrectGemmaTwentySeven}{27}
\newcommand{\tcmFollowCorrectMedGemma}{9}
\newcommand{\tcmFollowCorrectScout}{24}
\newcommand{\tcmFollowIncorrectGemmaTwelve}{14}
\newcommand{\tcmFollowIncorrectGemmaTwentySeven}{14}
\newcommand{\tcmFollowIncorrectMedGemma}{8}
\newcommand{\tcmFollowIncorrectScout}{17}
\newcommand{\tcmN}{52}
\newcommand{\tcmNonCont}{17}
\newcommand{\tcmNonN}{17}
\newcommand{\tcmPdEsc}{35}
\newcommand{\tcmPdN}{35}
\newcommand{\vfourBpN}{94}
\newcommand{\vfourBpRatio}{1.4}
\newcommand{\vfourBpRho}{0.57}
\newcommand{\vfourConcordant}{130}
\newcommand{\vfourConcordantPct}{49}
\newcommand{\vfourExpCRN}{15}
\newcommand{\vfourExpCRSame}{12}
\newcommand{\vfourExpPDN}{161}
\newcommand{\vfourExpPDSame}{110}
\newcommand{\vfourExpPRN}{16}
\newcommand{\vfourExpPRSame}{2}
\newcommand{\vfourExpSDN}{76}
\newcommand{\vfourExpSDSame}{6}
\newcommand{\vfourFollowups}{268}
\newcommand{\vfourMajorityAIA}{25.8}
\newcommand{\vfourMajorityDSCR}{56.5}
\newcommand{\vfourMajorityLIL}{31.4}
\newcommand{\vfourMajorityTCM}{43.5}
\newcommand{\vfourNAIA}{62}
\newcommand{\vfourNDSCR}{62}
\newcommand{\vfourNLIL}{51}
\newcommand{\vfourNPJRF}{61}
\newcommand{\vfourNTCM}{62}
\newcommand{\vfourPatients}{62}
\newcommand{\vfourPending}{1}
\newcommand{\vfourPjrfAucLR}{0.547}
\newcommand{\vfourPjrfBase}{55.7}
\newcommand{\vfourPjrfBrierBase}{0.255}
\newcommand{\vfourPjrfBrierLR}{0.264}
\newcommand{\vfourPjrfN}{61}
\newcommand{\vfourSelectedAgree}{27}
\newcommand{\vfourTcmConfirm}{19}
\newcommand{\vfourTcmContinue}{27}
\newcommand{\vfourTcmEscalate}{16}
\newcommand{\vfourTcmLookupConf}{74.2}
\newcommand{\vfourTcmLookupEsc}{69.4}
\newcommand{\vfourTcmPdItems}{35}
\newcommand{\vfourUndetermined}{37}
\newcommand{\vthreeDiffGemmaTwelve}{$-$2.7}
\newcommand{\vthreeDiffGemmaTwentySeven}{+3.1}
\newcommand{\vthreeDiffMedGemma}{$-$1.9}
\newcommand{\vthreeDiffScout}{$-$3.1}
\newcommand{\vthreeDscrAccGemmaTwelve}{55.8}
\newcommand{\vthreeDscrAccGemmaTwentySeven}{65.4}
\newcommand{\vthreeDscrAccMedGemma}{15.4}
\newcommand{\vthreeDscrAccScout}{57.7}
\newcommand{\vthreeDscrCtxHurtGemmaTwelve}{$-$15.4}
\newcommand{\vthreeDscrCtxHurtGemmaTwentySeven}{$-$17.3}
\newcommand{\vthreeDscrCtxHurtMedGemma}{$-$3.8}
\newcommand{\vthreeDscrCtxHurtPGemmaTwelve}{0.021}
\newcommand{\vthreeDscrCtxHurtPGemmaTwentySeven}{0.012}
\newcommand{\vthreeDscrCtxHurtPMedGemma}{0.625}
\newcommand{\vthreeDscrCtxHurtPScout}{0.039}
\newcommand{\vthreeDscrCtxHurtScout}{$-$13.5}
\newcommand{\vthreeDscrPGemmaTwelve}{0.97}
\newcommand{\vthreeDscrPGemmaTwentySeven}{0.68}
\newcommand{\vthreeDscrPMedGemma}{1.00}
\newcommand{\vthreeDscrPScout}{0.95}
\newcommand{\vthreeTcmDscrOnlyGemmaTwelve}{+15.4}
\newcommand{\vthreeTcmDscrOnlyGemmaTwentySeven}{+26.9}
\newcommand{\vthreeTcmDscrOnlyMedGemma}{+1.9}
\newcommand{\vthreeTcmDscrOnlyScout}{+26.9}
\newcommand{\vthreeTcmGapGemmaTwelve}{34.6}
\newcommand{\vthreeTcmGapGemmaTwentySeven}{38.5}
\newcommand{\vthreeTcmGapMedGemma}{17.3}
\newcommand{\vthreeTcmGapPGemmaTwelve}{0.000}
\newcommand{\vthreeTcmGapPGemmaTwentySeven}{0.000}
\newcommand{\vthreeTcmGapPMedGemma}{0.004}
\newcommand{\vthreeTcmGapPScout}{0.004}
\newcommand{\vthreeTcmGapScout}{17.3}
\newcommand{\vthreeTcmNoDscrGemmaTwelve}{+1.9}
\newcommand{\vthreeTcmNoDscrGemmaTwentySeven}{+5.8}
\newcommand{\vthreeTcmNoDscrMedGemma}{$-$3.8}
\newcommand{\vthreeTcmNoDscrScout}{+3.8}


% If the title and author information does not fit in the area allocated, uncomment the following
%
%\setlength\titlebox{<dim>}
%
% and set <dim> to something 5cm or larger.

% \title{Right Answer, No Image: Shortcut Learning in Brain-MRI Question Answering}   % earlier title; claims a mechanism the results do not show
\title{Do Medical MLLMs Read the Scan? Shortcut Controls on Brain-MRI Question Answering}

% Author information can be set in various styles:
% For several authors from the same institution:
% \author{Author 1 \and ... \and Author n \\
%         Address line \\ ... \\ Address line}
% if the names do not fit well on one line use
%         Author 1 \\ {\bf Author 2} \\ ... \\ {\bf Author n} \\
% For authors from different institutions:
% \author{Author 1 \\ Address line \\  ... \\ Address line
%         \And  ... \And
%         Author n \\ Address line \\ ... \\ Address line}
% To start a separate ``row'' of authors use \AND, as in
% \author{Author 1 \\ Address line \\  ... \\ Address line
%         \AND
%         Author 2 \\ Address line \\ ... \\ Address line \And
%         Author 3 \\ Address line \\ ... \\ Address line}

\author{First Author \\
  Affiliation / Address line 1 \\
  Affiliation / Address line 2 \\
  Affiliation / Address line 3 \\
  \texttt{email@domain} \\\And
  Song Wang \\
  Affiliation / Address line 1 \\
  Affiliation / Address line 2 \\
  Affiliation / Address line 3 \\
  \texttt{email@domain} \\}

\begin{document}
\maketitle
\begin{abstract}
Multimodal LLMs are scored on medical image questions by final-answer accuracy, which cannot show whether an answer came from the image, the question text, or a shortcut, and removing the image says nothing about image use when the question cannot be answered from it. We build five-phase question chains from the LUMIERE glioblastoma cohort and first audit an expert-labeled version: only \vfourConcordantPct{}\% of expert RANO ratings can be reproduced from the enhancing lesion on the displayed slices, so image-removal results on those items cannot separate shortcut use from unanswerable questions. We rebuild the chains so that each answer key is a stated function of what the model sees: which sequence is shown, where the lesion lies on the displayed slice, the enhancement-based response category, a one-year mortality forecast scored by the Brier score, and a management rule that depends on both the response category and the time since chemoradiotherapy. Controls are fixed before the runs: a positive control that only the image can answer, matched donor-image swaps in which the donor's answer is always available, fixed-context comparisons, and fact-flip tests of the management rule, applied to four open-weight models and one positive-control model. The control model reads the sequence and the lesion side from the image (\vfourEoneOwnAIAGemini{}\% and \vfourEoneOwnLILGemini{}\%, against \vfourEoneTextAIAGemini{}\% and \vfourEoneTextLILGemini{}\% without it), so those items are answerable from the slices; the open models mostly do not (\vfourEoneOwnAIAGemmaTwelve{} to \vfourEoneOwnAIAScout{}\% on sequence, chance 25\%), and none exceeds a category lookup on the management rule. We find no evidence of a learned shortcut as such, only of non-use of the image. Items and keys are not clinician-validated.
\end{abstract}

\section{Introduction}
\label{sec:introduction}

Medical multimodal LLMs are usually scored on the letter they pick, which cannot separate a model that read the image and reasoned validly from one that used cues in the question or options, or guessed. Such shortcuts are well documented outside medicine: hypotheses alone beat chance on natural language inference \citep{gururangan2018annotation}, models adopt syntactic heuristics that fail on adversarial sets \citep{mccoy2019right}, and stated chain-of-thought often does not reflect what produced the answer \citep{turpin2023language}. In chest-radiograph question answering, models often rely on linguistic patterns or attend to irrelevant image regions, and textual and visual perturbation protocols exist for detecting this \citep{nguyen2025localizing}. A control that removes or swaps the image is only informative when the question can be answered from the image: if the displayed slice does not contain the evidence for the key, performance will not change for a model that reads images perfectly.

We ask whether accuracy on brain-MRI questions reflects reading the scan or a shortcut, and we treat answerability as a property to be built and audited, not assumed. We use the five clinical phases of OmniBrainBench \citep{peng2026omnibrainbench}, anatomical and imaging assessment (AIA), lesion identification and localization (LIL), diagnostic synthesis and causal reasoning (DSCR), prognostic judgment and risk forecasting (PJRF), and therapeutic cycle management (TCM), arranged as a same-patient chain on the LUMIERE glioblastoma cohort \citep{suter2022lumiere}. Our first construction drafted items with an LLM from expert RANO ratings and cohort records. Auditing it shows that the keys are not functions of the model-visible inputs (Section~\ref{sec:audit}), so we rebuilt it (Section~\ref{sec:v4}) and fixed the controls before running them (Section~\ref{sec:controls}).

The paper makes three contributions:
\begin{itemize}[leftmargin=*, itemsep=2pt, topsep=2pt]
    \item An answerability audit of an expert-labeled brain-MRI chain: an enhancement-based rule applied to the displayed slices reproduces the expert RANO category for only \vfourConcordantPct{}\% of \vfourFollowups{} follow-ups, and the released benchmark's groupings cannot support same-patient chains at all (2 verifiable cross-phase links in 6,823 questions).
    \item A chain in which every key is a stated function of the model-visible inputs, and a control suite fixed in advance: a positive control that only the image can answer, matched donor-image swaps with the donor's answer always available, fixed-context comparisons, and fact-flip tests that separate use of the diagnosis label from use of other supplied facts.
    \item Results on four open-weight models and one positive-control model (Section~\ref{sec:results}): the control reads the image on the sequence and lesion-side questions, the open models mostly do not, and the diagnosis-category question cannot separate image use from a class prior even for the control. The KAB vocabulary, adaptation and gating metrics we proposed earlier are kept only as exploratory instruments (Appendix~\ref{sec:kab}).
\end{itemize}

\section{Related Work}
\label{sec:related-work}

\paragraph{Shortcuts and faithfulness.} A model can reach a label without doing the intended task: hypothesis-only baselines beat the majority class on many NLI datasets \citep{poliak2018hypothesis, gururangan2018annotation}, models trained on such data adopt heuristics that fail on HANS \citep{mccoy2019right}, and biasing features in the input change chain-of-thought answers without appearing in the explanation \citep{turpin2023language}.

\paragraph{Answering without the image.} In visual question answering, language priors let models answer without the image, and \citet{goyal2017making} rebalanced the data so each question has two images with different answers. \citet{chen2024right} find that visual content is unnecessary for many samples in popular multimodal benchmarks. In medicine, \citet{yan2025worse} show that state-of-the-art models score below random guessing on specialized diagnostic questions when ground-truth questions are paired with adversarial counterparts, \citet{jin2024hidden} find flawed rationales in 35.5\% of GPT-4V's correct answers, and \citet{nguyen2025localizing} introduce textual and visual perturbation tests for chest-radiograph questions. Their items are single questions; we apply the same logic to a same-patient chain and add a positive control, matched swaps, and interventions on the stated clinical facts. UniMod \citep{gu2026unimod} addresses the same failure at training time for image-plus-text diagnosis: it supervises image-only, text-only and fused predictions separately, and shows that a plainly fine-tuned fused model loses little when the image is removed but much more when the text is removed. We instead evaluate existing question-answering models and test whether the image is used at all, with a positive control that only the image can answer.

\paragraph{Brain benchmarks and same-patient continuity.} OmniBrainBench \citep{peng2026omnibrainbench} is a brain-imaging benchmark organized around a five-phase clinical workflow (15 modalities; 6,823 closed-ended questions). Its supplement describes a patient-level extension on the ADNI Alzheimer's cohort; the authors' repository provides label-generation code but no question data, and we did not evaluate it. LUMIERE \citep{suter2022lumiere} and MU-Glioma-Post \citep{muglioma2025post} provide real longitudinal glioma data with treatment and outcomes but no question layer.

\section{Building Answerable Chains from LUMIERE}
\label{sec:lumiere}

\paragraph{Why not OmniBrainBench.} A chain needs one patient's phases. We checked whether any image is shared across questions of different phases within a grouped case: of 6,823 closed-ended questions only 2 verifiable cross-phase patient links exist (both in VQA\_RAD, one AIA and one LIL question each), and a second method based on filename identifiers agreed; the counts are lower bounds (Appendix~\ref{sec:app-compat}). Each sub-corpus was built for one task, so a patient is essentially asked one phase-type of question. This is an incompatibility between the released groupings and our evaluation, not a defect in a promise its authors made.

\paragraph{LUMIERE.} LUMIERE \citep{suter2022lumiere} has 91 glioblastoma patients with longitudinal MRI (T1, contrast-enhanced T1, T2, FLAIR), automated segmentations (DeepBraTumIA and HD-GLIO-AUTO; \citealp{suter2022lumiere, kickingereder2019automated}), RANO ratings by one neuroradiologist \citep{wen2010updated}, and overall survival for most patients. Segmentations are automated and were not manually verified; the ratings are expert-assigned.

\subsection{Audit: are the expert labels answerable from the displayed slices?}
\label{sec:audit}

Our first item set (\texttt{v3}, Appendix~\ref{sec:app-v3}) let an LLM draft each phase's question from the expert RANO rating, the volume change between two segmentations, and cohort records. Three properties make its keys unanswerable from what the model sees. (i)~LIL asked for whole-lesion volume change, which one 2D slice per timepoint does not determine. (ii)~AIA had one fixed answer for every patient. (iii)~DSCR used the expert rating, which need not follow from the enhancing lesion. We measured (iii): on the axial slice with the largest enhancing area at each scan, we computed the RANO bidimensional product (longest diameter times longest perpendicular diameter) of the largest enhancing component and applied the enhancement criteria (progression: at least 25\% above the smallest earlier measurement, the nadir, or a new measurable lesion; response against the post-operative baseline; measurable means at least 10 by 10 mm). Of \vfourFollowups{} first-line follow-ups (one per rated, imaged timepoint before any second surgery), \vfourUndetermined{} are undetermined because no measurable enhancing disease exists on the baseline or nadir scan, and the rule reproduces the expert category for \vfourConcordant{} (\vfourConcordantPct{}\%): progressive disease \vfourExpPDSame{} of \vfourExpPDN{}, complete response \vfourExpCRSame{} of \vfourExpCRN{}, stable disease \vfourExpSDSame{} of \vfourExpSDN{}, partial response \vfourExpPRSame{} of \vfourExpPRN{}. The remaining ratings rest on evidence outside contrast enhancement on the displayed slices, such as T2/FLAIR progression, the choice of reference scan, clinical status and steroid dose. LUMIERE documents that one neuroradiologist rated the studies but not what information they used, so we do not say which. Our automated product correlates only moderately with the diameters the rater recorded for the target lesion (Spearman \vfourBpRho{} over \vfourBpN{} follow-ups; ours is a median \vfourBpRatio{} times larger), which is a limit on the segmentation-based keys. Agreement is a proxy: it does not show that the rule is right where the two disagree.

\subsection{The \texttt{v4} item set}
\label{sec:v4}

Each key is a stated function of the inputs the model receives (Table~\ref{tab:phases}). All random choices use fixed seeds and never look at model output.

\begin{table*}[t]
\centering
\footnotesize
\setlength{\tabcolsep}{3pt}
\begin{tabular}{p{0.7cm}p{4.6cm}p{4.4cm}p{4.4cm}r}
\toprule
Phase & Question & Model-visible inputs & Key & $n$ \\
\midrule
AIA & Which MRI sequence is shown? & one axial slice (T1, T1 after contrast, T2 or FLAIR) & the sequence rendered (balanced) & \vfourNAIA{} \\
LIL & Hemisphere and anterior/posterior half of the main tumor region & follow-up contrast-enhanced T1 slice at the level of the largest tumor region & position of the largest tumor-core component on that slice, relative to the midline and the brain's front-to-back midpoint & \vfourNLIL{} \\
DSCR & RANO enhancement category at follow-up & baseline, nadir and follow-up slices at the level of the largest enhancing lesion; the RANO thresholds are stated & category from the bidimensional products on those slices & \vfourNDSCR{} \\
PJRF & Probability of death within 52 weeks of this scan & age, sex, MGMT, IDH, extent of resection, weeks since surgery; alive at the scan & recorded outcome; scored by Brier score, not accuracy & \vfourNPJRF{} \\
TCM & Next management step & the same clinical facts plus weeks since chemoradiotherapy; the category comes from the chain & rule: non-PD continue; PD within 12 weeks of radiotherapy repeat MRI first; PD later change therapy & \vfourNTCM{} \\
\bottomrule
\end{tabular}
\caption{The \texttt{v4} phases. Stems for AIA, LIL and DSCR contain no patient-specific fact and every option set contains every possible answer, so a swapped image never conflicts with the text and the donor's answer is always available.}
\label{tab:phases}
\end{table*}

\paragraph{Cohort.} Of the patients with a first-line follow-up (rated, imaged, before any second surgery, with an imaged post-operative reference scan), we select one follow-up per patient in strata defined by the TCM rule class (continue, repeat MRI first, change therapy) and by the DSCR category, with quotas fixed in advance and a fixed seed, giving \vfourPatients{} patients. Selection used tabular records and segmentation masks only. Progressive-disease follow-ups whose timing falls within 2 weeks of the 12-week window are left out.

\paragraph{Images.} Each image is an atlas-space axial slice, skull-stripped by the dataset authors with HD-BET \citep{isensee2019automated}, at 4$\times$ magnification with a 20\,mm scale bar, in radiological orientation (checked against the affine). Slice levels are those on which the key is computed.

\paragraph{AIA, LIL, DSCR.} AIA is the positive control: its key is the sequence rendered, uniform over four sequences, so text alone can reach only chance. LIL's key is the hemisphere and anterior/posterior half of the largest tumor-core component on the displayed slice; lesions within 10\,mm of a midline, smaller than 100\,mm$^2$ or crossing the midline are not asked about. DSCR's stem states the RANO enhancement criteria, and its key is the category computed from the bidimensional products. The expert rating is recorded for each item but is not the key.

\paragraph{PJRF.} The prediction time is the follow-up scan, the available information is what the stem states plus the model's earlier assessment, and the outcome is death within 52 weeks of the scan, from the LUMIERE survival record (which we assume shares its origin with the week numbering; censoring is not documented). The model picks one of four probability bins (midpoints .10, .30, .50, .80) and is scored with the Brier score, not by whether it picked a ``correct'' bin. The recorded outcome is death in \vfourPjrfBase{}\% of \vfourPjrfN{} patients; a leave-one-out logistic model on the stated facts reaches Brier \vfourPjrfBrierLR{} against \vfourPjrfBrierBase{} for the constant base rate (AUC \vfourPjrfAucLR{}), so these facts support little beyond the base rate at this sample size.

\paragraph{TCM.} LUMIERE records no treatment after chemoradiotherapy, so a documented treatment is not available and the key is a rule-based recommendation, not what was done: non-progressive disease continues adjuvant temozolomide; progressive disease within 12 weeks of completing radiotherapy is confirmed with repeat MRI before any change; later progressive disease prompts a change of therapy. The 12-week window is from the RANO criteria \citep{wen2010updated}, which also allow progression inside the window when new enhancement lies outside the radiation field; we cannot assess the field, so the rule is a simplification. Radiotherapy dates are not in LUMIERE, so the stem states the weeks since chemoradiotherapy under an assumed standard schedule (ending 10 weeks after surgery). The same diagnosis therefore maps to different actions depending on the stated timing (\vfourTcmConfirm{} repeat-MRI and \vfourTcmEscalate{} change-therapy items among \vfourTcmPdItems{} progressive-disease items), which lets a lookup of the category alone be told apart from use of the timing fact.

\section{Controls and Pre-specified Analysis}
\label{sec:controls}

The analysis plan (\texttt{paper/preregistration\_v4.md}) was written before any \texttt{v4} model run; its hash is recorded in the project log, and the file is not externally timestamped. Models are the four open-weight models of the earlier study (MedGemma-4B, Gemma-3-12B, Gemma-3-27B, Llama-4-Scout; 4-bit, greedy decoding, 800 tokens; unparseable responses score incorrect). Each condition asks one question with a manufactured upstream context, so a phase's effect is not mixed with the model's own earlier answers.

\begin{itemize}[leftmargin=*, itemsep=1pt, topsep=2pt]
    \item \textbf{E1, image effect} (AIA, LIL, DSCR): accuracy with the image minus text only, paired by patient, with a patient-bootstrap interval. The margin is $\pm10$ percentage points, fixed in advance; ``within margin'' requires the whole interval inside it, and the interval half-width is about $1.96\sqrt{d/n}$ for a share $d$ of items whose correctness differs, so at $n=62$ the margin is reachable only if $d\le14\%$. Holm correction over the 12 model-phase tests; other contrasts are descriptive and unadjusted, with exact McNemar tests \citep{mcnemar1947note} and Wilson intervals \citep{wilson1927probable} where a proportion is reported. \textbf{AIA is the positive control:} for a model whose AIA interval is not above zero, the test has not been shown to detect image use, and E1 on LIL and DSCR supports no conclusion about it.
    \item \textbf{E2, donor tracking}: swap the image for a donor's with a different key (one donor per patient and phase, same number of images); $G=1[\text{swap}=\text{donor key}]-1[\text{text only}=\text{donor key}]$, so the no-image prior is subtracted.
    \item \textbf{E3, use of stated facts} (TCM): with the true context, move the stated timing across the 12-week window (\emph{flip\_window}) or replace the stated category by one that maps to another action class (\emph{flip\_label}), and report the share of answers that move to the rule's new answer. A category lookup alone cannot: it reaches at most \vfourTcmLookupConf{}\% (\vfourTcmLookupEsc{}\% if it always escalates), the rule 100\%.
    \item \textbf{E4, upstream context}: true, randomly wrong (each earlier answer replaced by a seeded random incorrect option) or absent context.
    \item \textbf{E5, PJRF}: Brier score and skill against the leave-one-out base rate, and AUC.
\end{itemize}

The ordinary five-phase chain (own answers as context) is run with and without images as a secondary analysis; each arm regenerates its upstream answers, so later-phase differences are total effects through the chain.

\paragraph{Control model.} Two of the four open models did not pass the positive control (Section~\ref{sec:results}), so before running any further model we added one control model to the plan (recorded in its change log): Gemini-3.6-Flash (\texttt{gemini-3.6-flash}, Google API, run 2026-09-25). It runs E1 and E2 only, with the same prompts, temperature 0 and decision rules. Two differences are forced by the API: the output cap is 16,384 tokens, because the model's hidden reasoning tokens count against it, and the route is the provider's OpenAI-compatible endpoint. The plan's Gemini-2.5-Pro was no longer offered to new accounts, and a Pro preview model was kept as the fallback should the Flash model fail the control. The control is reported apart from the four-model Holm family, empty API responses are counted and rerun (there were \vfourCtrlEmpty{} among \vfourCtrlRecords{} records), and, being an API model, it can change without notice. A second addition, also recorded before it ran, is a precision control: Gemma-3-27B is rerun in bf16 (Google's \texttt{gemma-3-27b-it} weights, vLLM, one GPU) on the same three conditions and phases, reported apart from the Holm family. Its reading rule was fixed in advance: an AIA interval above zero and clearly above the 4-bit model's would mean quantization explains part of the gap, and a result near chance that it does not. The serving stack differs from the Ollama runs, so precision and stack are confounded.

\section{Results}
\label{sec:results}

\subsection{Reference values without a model}

Table~\ref{tab:baselines} gives what constants and lookups reach. AIA is uniform over four answers, DSCR's most common category is progressive disease for \vfourMajorityDSCR{}\% of items, and a category-only TCM lookup reaches at most \vfourTcmLookupConf{}\%.

\begin{table}[t]
\centering
\footnotesize
\setlength{\tabcolsep}{3pt}
\begin{tabular}{lrrl}
\toprule
Phase & $n$ & Chance & Most common key \\
\midrule
AIA & \vfourNAIA{} & 25.0 & \vfourMajorityAIA{}\% \\
LIL & \vfourNLIL{} & 25.0 & \vfourMajorityLIL{}\% \\
DSCR & \vfourNDSCR{} & 25.0 & \vfourMajorityDSCR{}\% \\
TCM & \vfourNTCM{} & 25.0 & \vfourMajorityTCM{}\% (lookup $\le$ \vfourTcmLookupConf{}\%) \\
PJRF & \vfourNPJRF{} & (Brier) & base rate \vfourPjrfBase{}\% \\
\bottomrule
\end{tabular}
\caption{Non-model references on the \texttt{v4} items: chance, share of the most common key, the ceiling of a category-only TCM lookup, and the PJRF base rate.}
\label{tab:baselines}
\end{table}

\subsection{Image controls and stated facts}

\paragraph{The control model reads the image.} Gemini-3.6-Flash identifies the sequence for \vfourEoneOwnAIAGemini{}\% of patients with the image and \vfourEoneTextAIAGemini{}\% without it (difference \vfourEoneDiffAIAGemini{} points, 95\% CI [\vfourEoneLoAIAGemini{}, \vfourEoneHiAIAGemini{}]), and it places the lesion for \vfourEoneOwnLILGemini{}\% against \vfourEoneTextLILGemini{}\%. With the image swapped for a donor's it gives the donor's answer for \vfourEtwoSwapAIAGemini{}\% (AIA) and \vfourEtwoSwapLILGemini{}\% (LIL) of patients (Table~\ref{tab:ctrl}). These items are therefore answerable from the displayed slices, and the design detects image use when a model has it; that also means a weak result on them for another model is a property of that model. DSCR is different. The control scores \vfourEoneOwnDSCRGemini{}\% with the image and \vfourEoneTextDSCRGemini{}\% without ($\Delta$ \vfourEoneDiffDSCRGemini{} [\vfourEoneLoDSCRGemini{}, \vfourEoneHiDSCRGemini{}]), and its text-only accuracy equals the share of the most common key (\vfourMajorityDSCR{}\%: it answers progressive disease for every patient). It does use the image, since a swap changes \vfourEtwoChangesDSCRGemini{}\% of its answers (gain \vfourEtwoGainDSCRGemini{} [\vfourEtwoLoDSCRGemini{}, \vfourEtwoHiDSCRGemini{}]), but the image adds little to accuracy beyond the class prior. A null DSCR effect for another model therefore cannot be read as a failure to use the image.

\paragraph{Open models: image effect (E1).} Sequence identification is \vfourEoneOwnAIAGemmaTwelve{} to \vfourEoneOwnAIAScout{}\% (chance 25\%; Table~\ref{tab:e1}). By the pre-registered rule the positive control is passed by MedGemma-4B (\vfourEoneDiffAIAMedGemma{} [\vfourEoneLoAIAMedGemma{}, \vfourEoneHiAIAMedGemma{}], \vfourEoneOnlyOwnAIAMedGemma{} items right only with the image, \vfourEoneOnlyTextAIAMedGemma{} only without) and Llama-4-Scout (\vfourEoneDiffAIAScout{} [\vfourEoneLoAIAScout{}, \vfourEoneHiAIAScout{}]), whose gains are small beside the control's, and not by Gemma-3-12B (\vfourEoneDiffAIAGemmaTwelve{} [\vfourEoneLoAIAGemmaTwelve{}, \vfourEoneHiAIAGemmaTwelve{}]) or Gemma-3-27B (\vfourEoneDiffAIAGemmaTwentySeven{} [\vfourEoneLoAIAGemmaTwentySeven{}, \vfourEoneHiAIAGemmaTwentySeven{}]); for those two the plan permits no claim about LIL or DSCR image use. No E1 test survives Holm correction (smallest adjusted $p=\vfourEoneHolmAIAScout{}$). On LIL every interval includes zero, from \vfourEoneDiffLILMedGemma{} (MedGemma-4B) to \vfourEoneDiffLILScout{} (Llama-4-Scout); on DSCR Gemma-3-12B shows \vfourEoneDiffDSCRGemmaTwelve{} [\vfourEoneLoDSCRGemmaTwelve{}, \vfourEoneHiDSCRGemmaTwelve{}], which by the rule cannot be interpreted, and MedGemma-4B and Gemma-3-27B are within the margin.

\paragraph{Precision control.} Rerun in bf16, Gemma-3-27B identifies the sequence for \vfourEoneOwnAIAGemmaTwentySevenBf{}\% of patients with the image and \vfourEoneTextAIAGemmaTwentySevenBf{}\% without it ($\Delta$ \vfourEoneDiffAIAGemmaTwentySevenBf{} [\vfourEoneLoAIAGemmaTwentySevenBf{}, \vfourEoneHiAIAGemmaTwentySevenBf{}]), against \vfourEoneOwnAIAGemmaTwentySeven{}\% and \vfourEoneTextAIAGemmaTwentySeven{}\% at 4-bit and \vfourEoneOwnAIAGemini{}\% for the control model. The interval includes zero and the accuracy stays near chance (25\%), so by the rule fixed in advance quantization does not explain the gap to the control; on LIL and DSCR it is also indistinguishable from text-only (\vfourEoneDiffLILGemmaTwentySevenBf{} and \vfourEoneDiffDSCRGemmaTwentySevenBf{}), and donor tracking on AIA is \vfourEtwoGainAIAGemmaTwentySevenBf{} [\vfourEtwoLoAIAGemmaTwentySevenBf{}, \vfourEtwoHiAIAGemmaTwentySevenBf{}]. One model was rerun, on a different serving stack, so this does not rule out an effect of quantization in the other open models.

\paragraph{Open models: donor tracking (E2).} Only Llama-4-Scout on LIL shows clear tracking (gain \vfourEtwoGainLILScout{} [\vfourEtwoLoLILScout{}, \vfourEtwoHiLILScout{}]); Llama-4-Scout and Gemma-3-27B also track on AIA (\vfourEtwoGainAIAScout{} and \vfourEtwoGainAIAGemmaTwentySeven{}), and the Gemma-3-27B LIL gain is negative (\vfourEtwoGainLILGemmaTwentySeven{} [\vfourEtwoLoLILGemmaTwentySeven{}, \vfourEtwoHiLILGemmaTwentySeven{}]). A swap changes only \vfourEtwoChangesDSCRGemmaTwentySeven{}\% of Gemma-3-27B's DSCR answers. Every other interval includes zero, and even the clear cases are far below the control's (\vfourEtwoGainAIAGemini{} and \vfourEtwoGainLILGemini{}).

\paragraph{Stated facts and context (E3, E4).} With the true context TCM accuracy is \vfourEthreeGoldMedGemma{}, \vfourEthreeGoldGemmaTwelve{}, \vfourEthreeGoldGemmaTwentySeven{} and \vfourEthreeGoldScout{}\% for MedGemma-4B, Gemma-3-12B, Gemma-3-27B and Llama-4-Scout, all below the category-only lookup (at most \vfourTcmLookupConf{}\%), so the models do not even exploit the stated category as far as a lookup would. When the stated timing crosses the 12-week window, the share of progressive-disease answers that move to the rule's new answer is \vfourEthreeTimingMedGemma{}\% for MedGemma-4B and \vfourEthreeTimingGemmaTwelve{}, \vfourEthreeTimingGemmaTwentySeven{} and \vfourEthreeTimingScout{}\% for the others; MedGemma-4B ignores the timing, and it answers the most common class almost regardless (continue items \vfourEthreeGoldContinueMedGemma{}\% correct, the other two classes \vfourEthreeGoldConfirmMedGemma{}\% and \vfourEthreeGoldEscalateMedGemma{}\%). Replacing the stated category moves \vfourEthreeLabelMedGemma{}, \vfourEthreeLabelGemmaTwelve{}, \vfourEthreeLabelGemmaTwentySeven{} and \vfourEthreeLabelScout{}\% of answers to the rule's answer. True context beats random wrong context clearly only for Gemma-3-12B on TCM (\vfourEfourGoldWrongDiffGemmaTwelve{} [\vfourEfourGoldWrongLoGemmaTwelve{}, \vfourEfourGoldWrongHiGemmaTwelve{}]). For LIL and DSCR the true context does not help in any model, and the differences are within a few points.

\paragraph{PJRF (E5).} No model forecasts clearly better than the base rate: skill is \vfourEfiveSkillMedGemma{}, \vfourEfiveSkillGemmaTwelve{}, \vfourEfiveSkillGemmaTwentySeven{} and \vfourEfiveSkillScout{}, as the plan anticipated from the stated facts alone. MedGemma-4B and Gemma-3-12B choose the 0.5 bin for \vfourEfiveMidMedGemma{}\% and \vfourEfiveMidGemmaTwelve{}\% of patients, so their Brier score is close to 0.25 by construction. MedGemma-4B also declines to answer the PJRF question without upstream context (all \vfourNPJRF{} responses), which is scored as a 0.5 forecast.

\paragraph{Chain (secondary).} With the model's own answers as context, the image-minus-text differences are mostly small and their intervals include zero; the largest are Llama-4-Scout on AIA (\vfourEsixDiffAIAScout{} [\vfourEsixLoAIAScout{}, \vfourEsixHiAIAScout{}]) and TCM (\vfourEsixDiffTCMScout{} [\vfourEsixLoTCMScout{}, \vfourEsixHiTCMScout{}]). These are not corrected for multiplicity.

\paragraph{What this shows and does not show.} On sequence and lesion side, a control model uses the image almost completely while the open models mostly do not, and because the control answers these items the difference is not a matter of unanswerable questions. On the management rule no model beats a category lookup, and one ignores the stated timing. This is evidence of non-use of the image and of the stated facts, not of a learned shortcut as such: the open models are smaller and differ in training from the control, and either would produce the same gap. Quantization is separated only for Gemma-3-27B, whose bf16 rerun stays near chance; the other three open models ran at 4-bit only. Where the answer can be checked against a prior, as for DSCR, text-only accuracy equals the most common key for Gemma-3-27B and for the control. The DSCR question cannot separate image use from that prior even for the control, so we make no claim about it. None of this is validated by clinicians.


\begin{table*}[t]
\centering
\footnotesize
\setlength{\tabcolsep}{3pt}
\begin{tabular}{llrrrll}
\toprule
Model & Phase & $n$ & Image & Text-only & $\Delta$ [95\% CI] & Verdict \\
\midrule
MedGemma-4B & AIA / LIL / DSCR & -- & -- & -- & -- & -- \\
Gemma-3-12B & AIA / LIL / DSCR & -- & -- & -- & -- & -- \\
Gemma-3-27B & AIA / LIL / DSCR & -- & -- & -- & -- & -- \\
Llama-4-Scout & AIA / LIL / DSCR & -- & -- & -- & -- & -- \\
\bottomrule
\end{tabular}

\caption{E1: accuracy (\%) with the image and with text only, with no upstream context, and the paired difference in percentage points. AIA is the positive control.}
\label{tab:e1}
\end{table*}

\begin{table*}[t]
\centering
\footnotesize
\setlength{\tabcolsep}{3pt}
\begin{tabular}{llrrrrl}
\toprule
Model & Phase & $n$ & Own image & Swap: donor key & Text: donor key & Gain [95\% CI] \\
\midrule
MedGemma-4B & AIA / LIL / DSCR & -- & -- & -- & -- & -- \\
Gemma-3-12B & AIA / LIL / DSCR & -- & -- & -- & -- & -- \\
Gemma-3-27B & AIA / LIL / DSCR & -- & -- & -- & -- & -- \\
Llama-4-Scout & AIA / LIL / DSCR & -- & -- & -- & -- & -- \\
\bottomrule
\end{tabular}

\caption{E2: share of answers (\%) that equal the donor's key when the image is swapped, and when there is no image; the gain is the difference in percentage points.}
\label{tab:e2}
\end{table*}

\begin{table*}[t]
\centering
\footnotesize
\setlength{\tabcolsep}{3pt}
\begin{tabular}{lrrrlrl}
\toprule
Phase & $n$ & Image & Text-only & $\Delta$ [95\% CI] & Swap: donor key & Gain [95\% CI] \\
\midrule
AIA & 62 & 93.5 & 24.2 & +69.4 [+56.5, +80.6] & 93.5 & +72.6 [+61.3, +83.9] \\
LIL & 51 & 100.0 & 33.3 & +66.7 [+52.9, +78.4] & 100.0 & +74.5 [+62.7, +86.3] \\
DSCR & 62 & 61.3 & 56.5 & +4.8 [$-$6.5, +16.1] & 51.6 & +16.1 [+4.8, +27.4] \\
\bottomrule
\end{tabular}

\caption{Control model (Gemini-3.6-Flash): accuracy (\%) with the image and with text only, the paired difference, the share of swapped-image answers equal to the donor's key, and the donor-tracking gain (percentage points).}
\label{tab:ctrl}
\end{table*}

\begin{table}[t]
\centering
\footnotesize
\setlength{\tabcolsep}{3pt}
\begin{tabular}{lccccc}
\toprule
 & \multicolumn{3}{c}{TCM acc.\ (\%)} & \multicolumn{2}{c}{Follows rule (\%)} \\
\cmidrule(lr){2-4}\cmidrule(lr){5-6}
 & true & wrong & none & label & timing (PD) \\
\midrule
MedGemma-4B & -- & -- & -- & -- & -- \\
Gemma-3-12B & -- & -- & -- & -- & -- \\
Gemma-3-27B & -- & -- & -- & -- & -- \\
Llama-4-Scout & -- & -- & -- & -- & -- \\
\bottomrule
\end{tabular}

\caption{E3: TCM accuracy under true, random wrong and absent upstream context, and the share of answers that move to the rule's answer after the stated category is changed (label) or the stated timing crosses the 12-week window (timing, progressive-disease items).}
\label{tab:e3}
\end{table}

\begin{table}[t]
\centering
\footnotesize
\setlength{\tabcolsep}{3pt}
\begin{tabular}{lccccc}
\toprule
Model & $n$ & Brier & Base rate & Skill & AUC \\
\midrule
MedGemma-4B & -- & -- & -- & -- & -- \\
Gemma-3-12B & -- & -- & -- & -- & -- \\
Gemma-3-27B & -- & -- & -- & -- & -- \\
Llama-4-Scout & -- & -- & -- & -- & -- \\
\bottomrule
\end{tabular}

\caption{E5: PJRF Brier score with the true upstream context, against the leave-one-out constant base rate; skill $=1-$ Brier / base-rate Brier.}
\label{tab:e5}
\end{table}

\subsection{What the first item set showed}
\label{sec:v3-summary}

On \texttt{v3} (Appendix~\ref{sec:app-v3}; exploratory, with a margin chosen after seeing the intervals) removing the images changed accuracy by \vthreeDiffScout{}, \vthreeDiffGemmaTwentySeven{}, \vthreeDiffGemmaTwelve{} and \vthreeDiffMedGemma{} percentage points for Llama-4-Scout, Gemma-3-27B, Gemma-3-12B and MedGemma-4B. Given the audit above this is what one would see whether or not the models read the images, and we draw no shortcut conclusion from it. Two \texttt{v3} results are about the items' structure. TCM keys follow the DSCR category by construction (\tcmAgree{} of \tcmN{} keys map from the category to their action class), and accordingly correct upstream context beats forced-wrong context for all four models. DSCR is a class-imbalanced phase: progressive disease is the key for \dscrMajorityCount{} of \nPatVthree{} items (\dscrMajorityPct{}\%), above every model's DSCR accuracy. The motivation for the \texttt{v4} controls is that neither result can tell use of the category from use of the other stated facts.

\section{Conclusion}
\label{sec:conclusion}

Accuracy on an expert-labeled brain-MRI chain cannot show whether models read the scan: only \vfourConcordantPct{}\% of expert RANO ratings are reproducible from the enhancing lesion on the displayed slices, and the benchmark we started from cannot support same-patient chains. We rebuilt the chain on LUMIERE so that every key is a stated function of the model-visible inputs, and fixed a control suite in advance, with a positive control that only the image can answer, matched donor swaps, fixed-context comparisons and fact-flip tests of a management rule. A control model reads the sequence and the lesion side from the image (\vfourEoneOwnAIAGemini{}\% and \vfourEoneOwnLILGemini{}\%, against \vfourEoneTextAIAGemini{}\% and \vfourEoneTextLILGemini{}\% without it), so those items are answerable from the slices; four open-weight models mostly do not, only Llama-4-Scout tracks a swapped image on lesion side, and no model exceeds a category lookup on the management rule. The diagnosis-category question cannot separate image use from a class prior even for the control. These results show non-use of the image, not a learned shortcut as such: model scale and training are not separated from it, and 4-bit serving is separated only for Gemma-3-27B. Clinician review of the keys, and in particular of the management rule and the assumed radiotherapy schedule, is the outstanding step, and only one proprietary control model is run. The controls apply to any chain-structured medical benchmark whose keys are made answerable from its inputs.


\section*{Limitations}
\label{sec:limitations}

\textbf{Clinician review is outstanding.} No clinician has judged whether each question is answerable from the inputs shown or whether each key is right; ``answerable by construction'' means only that each key is a stated function of the inputs. The functions rest on automated segmentation (DeepBraTumIA), whose errors propagate to the LIL and DSCR keys. The enhancement rule ignores T2/FLAIR change, clinical status and steroid dose, abstains where no measurable enhancing disease exists, and reproduces the expert category for \vfourConcordantPct{}\% of follow-ups; agreement does not show that the rule is right where the two disagree, and the expert rating is recorded but is not the key.

\textbf{The management key is a rule, not a documented treatment.} LUMIERE records no treatment after chemoradiotherapy. The TCM key rests on an assumed radiotherapy schedule (ending 10 weeks after surgery) and on a simplified reading of the 12-week RANO window, which also allows progression inside the window when new enhancement lies outside the radiation field. The PJRF outcome assumes that survival time and scan weeks share the surgery-week origin, and censoring is not documented. With \vfourPjrfN{} patients the stated facts do not support a forecast better than the base rate (leave-one-out Brier \vfourPjrfBrierLR{} against \vfourPjrfBrierBase{}), so PJRF can only show that a model is not worse than a constant.

\textbf{Power and scope.} With 51 to 62 patients per phase, the pre-specified $\pm10$-point margin is reachable only when few items differ between arms, and several cells will be inconclusive. The positive control tests one skill, sequence identification; passing it does not show that a model can read lesions. The donor images come from different patients, so anatomy differs beyond the tested feature; stems in the text-only arm still refer to images. Patients are from one center and the cohort was stratified by design; the four open-weight models run at 4-bit through Ollama while the control is an API model of unknown size and precision, so the gap between them mixes scale, quantization and training and cannot be attributed to shortcut learning; a bf16 rerun of Gemma-3-27B (through vLLM, so precision is confounded with the serving stack) stays near chance on AIA, which argues against quantization as the main cause for that model but is not shown for the others. There is one control model, run once at temperature 0, and API models change without notice. Intervals treat patients as exchangeable.

\textbf{The first item set is exploratory.} Its items and distractors were LLM-drafted and are unreviewed; its equivalence margin was set after seeing the intervals; the donor-permutation, gating-causality, ablation and options-only analyses were added after seeing earlier results; 48 gating contrasts are uncorrected; the wrong-context condition used one fixed incorrect option; and the own-DSCR split of TCM accuracy is confounded by patient difficulty. KAB's ontology coverage is a string-matching metric that measures term validity, not correctness, and answer recoverability detects only inconsistency between reasoning and answer.

\section*{Ethics Statement}

This work evaluates the diagnostic reasoning of general-purpose and medical-domain language models on real, de-identified patient imaging data; it does not deploy, recommend, or endorse any model for clinical use, and no results in this paper should be interpreted as a clinical validation of any evaluated system. LUMIERE is used under its published non-commercial research license, and all patient data is accessed in its already de-identified, publicly released form; no new identifiable patient information is collected.

% Bibliography entries for the entire Anthology, followed by custom entries
%\bibliography{custom,anthology-overleaf-1,anthology-overleaf-2}

% Custom bibliography entries only
\bibliography{custom}

\appendix

\section{The First Item Set (\texttt{v3}) and Its Controls}
\label{sec:app-v3}

This appendix records the LLM-drafted item set of the earlier study and the controls run on it. It is exploratory: its keys are not functions of the model-visible inputs (Section~\ref{sec:audit}), and several analyses were added after seeing earlier results.

\subsection{Construction}
\label{sec:app-v3-construction}

LUMIERE \citep{suter2022lumiere} has 91 glioblastoma patients with automated multi-sequence segmentation and expert RANO ratings. Facts come from its shipped data (segmentation masks already in MNI152 space, precomputed volumes, survival and Stupp-protocol records); DSCR uses the expert ratings directly, which the displayed slices reproduce for only \vfourConcordantPct{}\% of follow-ups (Section~\ref{sec:audit}). Of 91 patients, 55 are eligible (at least 3 response-rated follow-ups and one timepoint with complete imaging); 54 were drafted, giving 269 items, and \nPatVthree{} remain in the final set. Claude 4.5 Sonnet drafted one item per phase with distractors from the facts. Each patient contributes one random eligible timepoint with both a rating and complete imaging (fixed seed); always taking the last rated timepoint would make nearly every DSCR answer progressive disease. DSCR labels over the \nPatVthree{} patients are \dscrVthreePD{} progressive disease, \dscrVthreeSD{} stable disease, \dscrVthreePR{} partial response and \dscrVthreeCR{} complete response. An automatic stem-leakage auditor flags any stem or option restating an answer of that or an upstream phase, with up to eight redraft retries; the final set has zero violations, which shows compliance with an auditor the drafting loop was optimized against, not the absence of leakage. Each timepoint is one axial slice (largest tumor area, contrast-enhanced T1, 1st to 99th percentile scaling), so LIL's whole-lesion volume change is not determined by what is shown. The development history (two sampling bugs, a leakier first set) is in \texttt{paper/update.md}.

\subsection{Setup}
\label{sec:setup}

Four open-weight models are served through Ollama at 4-bit (Q4\_K\_M) on one RTX PRO 6000 each: MedGemma-4B, Gemma-3-12B, Gemma-3-27B and Llama-4-Scout (109B mixture-of-experts, 17B active). All use one OpenAI-compatible interface, temperature 0 and 800 tokens; the schema asks for the answer before grounding and reasoning, so a truncated response still has one. A call is retried three times; unparseable JSON is salvaged by regex on the answer letter and otherwise scored incorrect, not excluded, so accuracy is per question. RadLex and the NCI Thesaurus \citep{sioutos2007nci} are used as-is.

\subsection{Controls}
\label{sec:grounding-test}

Recoverability and ontology coverage test internal consistency, not whether reasoning derives from the image, so we ran controls identically across the four models: a text-only ablation (chain context regenerated from each arm's own answers); an evidence fact-check of stated claims against the structured facts (availability in shown text, and record status; deterministic pattern matching, no paraphrase); counterfactual image substitution, with a donor-permutation null that redraws each patient's donor among patients with the opposite LIL direction (20,000 replicates); a repeated identical-input run; a gating-causality intervention, in which LIL, DSCR, PJRF and TCM are asked under true, fixed-wrong and absent upstream context, plus a TCM ablation that injects only the DSCR entry or every entry except it\label{sec:gating-causality}; and an options-only baseline. Contrasts are paired on the patient with bootstrap intervals and exact McNemar tests, uncorrected. A joint decision rule was recorded in the dated project log before the first jobs (2026-09-22; not externally timestamped, so a plan, not a preregistration); the $\pm8$pp margin was chosen after seeing the intervals, and the donor-permutation, gating, ablation and options-only analyses are exploratory.

\subsection{Results}
\label{sec:app-v3-results}

\begin{table}[t]
\centering
\footnotesize
\setlength{\tabcolsep}{3pt}
\begin{tabular}{lccc}
\toprule
Model & Image & Text-only & $\Delta$ [95\% CI] \\
\midrule
Llama-4-Scout & 48.8 & 51.9 & $-$3.1 [$-$7.3, +1.2] \\
Gemma-3-27B & 53.5 & 50.4 & +3.1 [$-$0.8, +7.3] \\
Gemma-3-12B & 48.1 & 50.8 & $-$2.7 [$-$7.3, +2.3] \\
MedGemma-4B & 34.2 & 36.2 & $-$1.9 [$-$5.8, +1.9] \\
\bottomrule
\end{tabular}

\caption{Non-gated accuracy (\%) with images versus text-only, and the paired difference (image $-$ text-only, points) with a patient-clustered 95\% bootstrap interval (52 patients).}
\label{tab:textonly}
\end{table}

\paragraph{Text-only ablation.} Three of four models score at least as high without the image (Table~\ref{tab:textonly}; Llama-4-Scout \vthreeDiffScout{}, Gemma-3-12B \vthreeDiffGemmaTwelve{}, MedGemma-4B \vthreeDiffMedGemma{} points; Gemma-3-27B \vthreeDiffGemmaTwentySeven{}, mostly AIA). Every interval includes zero and lies within $\pm8$pp. Models give the same answer with and without the image on 77 to 89\% of items, so the image changes answers without changing net accuracy. Repeating the own-image run gives identical answers on 260 of 260 questions for all four models, so every change below is caused by the input.

\paragraph{Evidence fact-check.} 26 to 55\% of answers contain no checkable claim. Among the rest (251 to 668 claims per model), 75 to 87\% were available in shown text and match the record, 10 to 18\% were available yet contradict it (for example a wrong location repeated from an earlier turn), 1 to 2\% were novel and match, and 2 to 6\% were novel and contradict. Few novel matching claims means little evidence that stated evidence contains facts obtainable only from the image; it cannot show the converse.

\paragraph{Counterfactual images.} Flip rates between own and donor image are 10 to 17\% for AIA (whose answer is a fixed string), 12 to 42\% for LIL and 6 to 27\% for DSCR. Among 95 flipped LIL and DSCR items, \permObserved{} (33.7\%) carry the donor's label against a permutation mean of \permMean{} (95\% range \permLo{} to \permHi{}; $p=\permP{}$), so which donor's image was shown does not matter beyond label base rates. The LIL part is degenerate by the pairing rule, and the donor's fully correct answer was among the options for 0 of 52 LIL patients, so the test is limited to label level.

\begin{table}[t]
\centering
\footnotesize
\setlength{\tabcolsep}{2pt}
\begin{tabular}{lcccc}
\toprule
 & \multicolumn{3}{c}{TCM} & DSCR \\
\cmidrule(lr){2-4}\cmidrule(lr){5-5}
Model & c / w / a & $\Delta_{\text{DSCR}}$ & $\Delta_{\text{no DSCR}}$ & c / w / a \\
\midrule
MedGemma-4B & 33/15/25 & +1.9 & $-$3.8 & 15/21/19 \\
Gemma-3-12B & 62/27/19 & +15.4$^{*}$ & +1.9 & 50/65/65 \\
Gemma-3-27B & 63/25/33 & +26.9$^{*}$ & +5.8 & 54/63/71 \\
Llama-4-Scout & 60/42/33 & +26.9$^{*}$ & +3.8 & 50/62/63 \\
\bottomrule
\end{tabular}

\caption{Accuracy (\%) under correct (c), forced-wrong (w) and absent (a) upstream context ($n=52$). $\Delta$ is correct minus wrong TCM accuracy when only the DSCR entry or every entry except DSCR is injected; $^{*}$ interval excludes zero. In the DSCR column the true context is AIA and LIL, never DSCR's own label. Uncorrected.}
\label{tab:gating}
\end{table}

\paragraph{Gating and TCM.} TCM depends on upstream context for all four models: correct beats forced-wrong context by \vthreeTcmGapMedGemma{}, \vthreeTcmGapGemmaTwelve{}, \vthreeTcmGapGemmaTwentySeven{} and \vthreeTcmGapScout{} points (Table~\ref{tab:gating}; $p\le.004$), and the DSCR entry alone reproduces the gap for three models. LIL and PJRF show no dependence. True context \emph{lowered} DSCR accuracy in three models (for example Gemma-3-27B \vthreeDscrCtxHurtGemmaTwentySeven{} points, $p=\vthreeDscrCtxHurtPGemmaTwentySeven{}$; uncorrected over 48 contrasts). The DSCR key predicts the TCM key's action class in \tcmAgree{} of \tcmN{} patients, so the dependency is built into the items, and TCM accuracy in the larger models reflects at least in part a lookup of the upstream category.

\paragraph{Options-only.} Options alone score 19 to 44\% on DSCR (below the \dscrMajorityPct{}\% majority rate), 21 to 25\% on PJRF and 12 to 19\% on TCM, so options are not the main shortcut there; on LIL (33 to 37\%) they match the same items with image and stem. AIA is the exception because its answer is the same string in every item. Always answering progressive disease scores \dscrMajorityPct{}\% on DSCR, above every model (\vthreeDscrAccGemmaTwentySeven{}, \vthreeDscrAccScout{}, \vthreeDscrAccGemmaTwelve{} and \vthreeDscrAccMedGemma{}\%), so no model shows RANO skill beyond class imbalance. MedGemma-4B's options-only responses were unparseable for 35 to 100\% of items in four phases, so those cells are not interpreted.

\paragraph{Reading.} On \texttt{v3}, removing images changed accuracy by at most 3.1 points with every interval including zero, swapped images did not move answers toward the donor beyond a permutation null, and TCM depends on the upstream category its keys encode. The flip rate is not low, so the plan's third condition is not met as written. We do not conclude that the images are ignored: the LIL, DSCR and AIA keys are not functions of the displayed slices, so these items cannot show image dependence, and no clinician has checked the keys.

\section{The KAB Framework (exploratory instruments)}
\label{sec:kab}
\label{sec:app-metrics}

KAB, the framework we proposed earlier, evaluates a case as a five-phase chain (AIA, LIL, DSCR, PJRF, TCM) and adds three measures. We keep it as an exploratory instrument and make no claim from it: recoverability is near ceiling, ontology coverage measures vocabulary membership, adaptation triggered on 1 to 4 of 52 patients per model, and gating is a scoring convention with no gated result reported. Each question gives the model the image, question, options and its own earlier answers (the chain context), and it returns an answer, a description of what it observed and its reasoning. \emph{Answer recoverability}: a second call sees only the reasoning and names the answer it supports, so a match shows the reasoning is textually sufficient, not that it caused the answer. \emph{Ontology term coverage}: the share of the response's knowledge-base terms in the phase-appropriate base (RadLex for AIA and LIL, the NCI Thesaurus for the rest); it measures term validity, not correctness. \emph{Feedback-adaptation}: when coverage falls below $\tau_g=0.3$ and concepts were missed, the phase is re-asked with the prior answer and the missed concepts, which does not separate a knowledge gap from an activation failure. \emph{Threshold gating}: a phase scores $0$ unless every phase it is assumed to depend on scored at least $\tau_c=0.5$; this is an imposed convention, tested instead by the upstream-context manipulation of Appendix~\ref{sec:gating-causality}.

\section{OmniBrainBench Compatibility Details}
\label{sec:app-compat}

A chain presupposes that a case's phases are one patient's trajectory. We tested this on OmniBrainBench \citep{peng2026omnibrainbench} by checking whether any image is shared across questions of different phases within a grouped case. In the 6,823-question closed-ended subset (15 source sub-corpora) only \textbf{2} verifiable cross-phase patient links exist, both in VQA\_RAD (one AIA and one LIL question each); a second method, stripping modality, view and slice suffixes from filenames, gave the same result. The release has no complete patient identifier, so these are lower bounds. The cause is structural: each sub-corpus was built for one clinical task, so a patient is essentially asked one phase-type of question, and reading \texttt{source\_file} as patient identity was our choice that stitches together different patients. This is an incompatibility with our evaluation, not a defect in a promise its authors made. The supplement's patient-level ADNI extension is not a brain-tumor cohort and, as of September 2026, has label-generation scripts only; we did not evaluate it. We found no replacement dataset with same-patient, multi-phase continuity and a question layer (Section~\ref{sec:related-work}), which motivates the LUMIERE construction.

\section{Evidence-Audit, Option-Only, and TCM-Lookup Details}
\label{sec:app-audit}

\paragraph{Evidence audit.} For each checkable claim in a response's grounding and reasoning (hemisphere or region statements, and measurements with units) we record \emph{availability} (already in the stem, options or upstream context) and \emph{record status}: \emph{supported} (a measurement within 2\% of a fact value, or a hemisphere present in the segmentation), \emph{contradicted} (a hemisphere opposite to the segmentation, assuming image left is patient right), \emph{unmatched} (a measurement matching no fact value, which may be a derived quantity) or \emph{unverifiable}. Keeping the two dimensions apart matters: an earlier single ``copied'' label hid a wrong claim repeated from an earlier turn. The audit does not interpret paraphrase and is not validated against expert annotation.

\paragraph{Option-only baselines} (\texttt{v3}, options shuffled, $n=52$). Always choosing the longest option scores 37\% on AIA, 25\% on LIL, 67\% on DSCR (the majority class, since progressive disease is typically the longest option), 19\% on PJRF and 6\% on TCM. The most common correct letter, fitted on the same items, scores 31 to 37\% on every phase. Except for AIA (a constant correct string) and DSCR (class imbalance), these sit near four-option chance.

\paragraph{TCM label lookup.} A fixed keyword rule (\texttt{tools/lumiere\_tcm\_lookup.py}) assigns each TCM option an action class from its first ten words. The DSCR key predicts the TCM key's class (progressive disease to escalate, other ratings to continue) in \tcmAgree{} of \tcmN{} patients (\tcmPdEsc{} of \tcmPdN{} and \tcmNonCont{} of \tcmNonN{}), but the class is unique among the four options in only \tcmClassUnique{} of \tcmN{}, so the class alone fixes the letter in a minority of items. Among the \tcmPdN{} progressive-disease patients, escalate-class TCM answers under correct, forced-wrong and absent DSCR context are \tcmFollowCorrectGemmaTwelve{}, \tcmFollowIncorrectGemmaTwelve{} and \tcmFollowAbsentGemmaTwelve{} (Gemma-3-12B); \tcmFollowCorrectGemmaTwentySeven{}, \tcmFollowIncorrectGemmaTwentySeven{} and \tcmFollowAbsentGemmaTwentySeven{} (Gemma-3-27B); \tcmFollowCorrectScout{}, \tcmFollowIncorrectScout{} and \tcmFollowAbsentScout{} (Llama-4-Scout); and \tcmFollowCorrectMedGemma{}, \tcmFollowIncorrectMedGemma{} and \tcmFollowAbsentMedGemma{} (MedGemma-4B).

\section{A Worked Patient Example}
\label{sec:worked-example}

Patient-002's \texttt{v3} chain, as answered by Gemma-3-27B (first case in the evaluation order, not selected). The patient is a 71-year-old female with glioblastoma (MGMT unmethylated), complete resection and Stupp chemoradiation; total lesion volume rises 27,937 to 57,570\,mm$^3$ (+106\%) between the post-operative and follow-up scans, the rating is progressive disease, and recorded survival is 48 weeks. The AIA answer is correct, but the model's grounding also says ``a mass lesion in the right temporal lobe'', which nothing shown mentions and which the record contradicts (the hemisphere check labels it a laterality flip). LIL repeats the claim and then gives the correct location and +106\%; under the two-dimension scoring of Appendix~\ref{sec:grounding-test} that response has 3 claims available in shown text that match the record and 2 (the repeated location) that contradict it, so the carried-forward error stays visible. DSCR (progressive disease) and TCM (second-line therapy) are correct. PJRF is the one wrong answer and shows a limit of the key: it is the bucket for the realized 48-week survival (about 12 months), the model chose about 9 months for the same adverse factors, and that forecast is plausible. Four of five answers are correct, but the evidence behind LIL and DSCR rests on a spatial claim introduced at AIA that the audit flags as contradicting the record, not on observation re-derived at each phase.

\end{document}
